\section{Coefficient comparison and copy-weight recovery}
\label{sec:sector_bnt_coeff_comparison}

% Source: copy-weight recovery, Cirac2016MPDO_arXiv Appendix MPV proof,
% lines 1184--1188.
The exact matching of Section~\ref{sec:sector_bnt_strong_match} produces, on each
matched pair, a coefficient identity. This section turns that identity into
equality of the per-copy weight multisets \cite[Appendix~A]{Cirac2016MPDO_arXiv}.
No analytic non-decay estimate enters: the fixed-length linear independence
isolates each sector coefficient exactly.

\subsection{Matched-sector weight multiset equality}

Given a matched pair $(j_0, \tilde k_0)$ and a nonzero scalar $\zeta$ with
$c_N^{(j_0)}(P) = \zeta^N c_N^{(\tilde k_0)}(Q)$ past a threshold $N_0$, this
section records two conclusions: the multiset equality of the per-copy weights
from power-sum recovery (Theorem~\ref{thm:power_sum_multiset}), and its upgrade
to an explicit per-copy permutation. Together they give the
$\mu_{j,q} = \nu_{j,q}\, e^{i\varphi_j}$ recovery of
\cite[Appendix~A]{Cirac2016MPDO_arXiv}.
% Source: matched normal sectors at Cirac2016MPDO_arXiv line 1182; equal-MPV
% coefficient comparison at lines 1184--1188; global-gauge construction at
% lines 1189--1192.
In the equal-MPV case, Chapter~\ref{ch:ft_proof} obtains this coefficient
identity by expanding both total MPVs in the matched BNT basis and applying
eventual linear independence
(Lemma~\ref{lem:sector_bnt_coeff_identity_via_matched_mpv_phase}). For merely
proportional total MPVs, a common length-dependent scalar remains in every
matched coefficient identity and the scalar-free comparison need not hold; see
Remark~\ref{rem:proportional_length_scalar_gap}.

\begin{lemma}[Matched-sector weight multiset equality]%
    \label{lem:sector_bnt_matched_sector_weight_multiset_eq}
    \lean{MPSTensor.matched_sector_weight_multiset_eq}
    \leanok
    \uses{def:sector_weight_data,
        thm:power_sum_multiset}
    Let $P, Q$ be sector decompositions with weights $\mu_{j,q}$ and
    $\nu_{k,q'}$, and fix indices $j_0$, $\tilde k_0$, and a nonzero scalar
    $\zeta \in \C \setminus \{0\}$. Suppose that for every $N > N_0$,
    \begin{align}
        c_N^{(j_0)}(P)
         & = \zeta^N \cdot c_N^{(\tilde k_0)}(Q).
        \label{eq:bnt_matched_sector_coeff_identity}
    \end{align}
    Then the per-sector multiplicities agree, $r_{j_0}(P) = r_{\tilde k_0}(Q)$,
    and the rescaled per-copy weight multisets coincide.
    % Source: Cirac2016MPDO_arXiv Appendix MPV proof, lines 1184--1188
    % (\mu_{j,q} = \nu_{j,q} e^{i\varphi_j}, matching multiplicities r_{a,j} = r_{b,j}).
    This is \cite[Appendix~A]{Cirac2016MPDO_arXiv} read on a single matched
    sector.
\end{lemma}

\begin{proof}\leanok
    \uses{thm:bounded_power_sum_multiset}
    Expand both sides of the coefficient identity
    (\ref{eq:bnt_matched_sector_coeff_identity}) into per-copy power sums
    $\sum_q \mu_{j_0, q}^N$ and $\sum_{q'} \nu_{\tilde k_0, q'}^N$, and
    distribute $\zeta^N$ through the right-hand sum to obtain, for every
    $N > N_0$,
    \begin{align}
        \sum_q \mu_{j_0, q}^N
         & = \sum_{q'} (\zeta \cdot \nu_{\tilde k_0, q'})^N.
        \notag
    \end{align}
    Extend to all positive exponents and apply the unequal-cardinality
    multiset-recovery lemma (Theorem~\ref{thm:bounded_power_sum_multiset}),
    using $\zeta \neq 0$ together with the weight non-vanishing assumption
    to ensure the input weights are nonzero. This yields the multiplicity
    match and the multiset equality.
\end{proof}

\begin{lemma}[Matched-sector weight-permutation extractor]%
    \label{lem:sector_bnt_matched_sector_weight_equiv}
    \lean{MPSTensor.matched_sector_weight_equiv}
    \leanok
    \uses{lem:sector_bnt_matched_sector_weight_multiset_eq}
    Under the hypotheses of
    Lemma~\ref{lem:sector_bnt_matched_sector_weight_multiset_eq}, the
    per-sector multiplicities agree and there is a permutation $\tau$ of
    the per-copy index set identifying the individual weights up to the
    gauge phase: for every $q$,
    $\nu_{\tilde k_0, \tau(q)} = \zeta^{-1} \cdot \mu_{j_0, q}$.
    % Source: Cirac2016MPDO_arXiv Appendix MPV proof, line 1188
    % (\mu_{j,q} = \nu_{j,q} e^{i\varphi_j}).
    This is the per-copy realisation of \cite[Appendix~A]{Cirac2016MPDO_arXiv}:
    the multiset equality from
    Lemma~\ref{lem:sector_bnt_matched_sector_weight_multiset_eq} is upgraded to
    a concrete indexing between the matched $P$- and $Q$-copies.
\end{lemma}

\begin{proof}\leanok
    Invoke Lemma~\ref{lem:sector_bnt_matched_sector_weight_multiset_eq} to
    obtain the cardinality match and the multiset equality.
    From this multiset-map equality extract a permutation matching the
    elements pointwise (proved by induction on the cardinality).
    Compose with the cardinality identification and cancel $\zeta$ via
    $\zeta \neq 0$ to convert the gauge factor on the right of the
    pointwise identity to $\zeta^{-1}$ on the left.
\end{proof}

\section{Strong existential matching by exact linear independence}
\label{sec:sector_bnt_strong_match}

The exact matching reads \cite[Appendix~A]{Cirac2016MPDO_arXiv} as a statement
at fixed length: equal MPV families force each basis sector of one form to pair
with a basis sector of the other. The non-vanishing of the sector coefficients
isolates the matched sector exactly, with no limit of a single-sector
projection.

\begin{theorem}[Exact single-sector matching]
    \label{thm:sector_bnt_exact_block_match}
    \lean{MPSTensor.exists_block_match_exact}
    \leanok
    \uses{def:sector_bnt_cf, def:same_mpv2_pos,
        thm:sector_bnt_proportional_exact_match}
    Let $P, Q$ be two BNT canonical forms with the same MPV family (at all
    positive lengths). For every sector $j$ of $P$ there exist a sector $k$ of
    $Q$ with matched bond dimension $\dim_P(j) = \dim_Q(k)$, a gauge-phase
    equivalence after identifying the matched bond dimensions, and a
    non-decaying cross-overlap. No per-sector unit-modulus hypothesis is
    required.
\end{theorem}

\begin{proof}\leanok
    \uses{thm:sector_bnt_proportional_exact_match}
    Equality at every positive length gives eventual nonzero proportionality
    with scalar $1$. Apply
    Theorem~\ref{thm:sector_bnt_proportional_exact_match} to the chosen sector.
\end{proof}

\begin{theorem}[Strong existential matching, full-basis form]
    \label{thm:sector_bnt_forall_unit_k_exists_j_nondecaying_overlap_of_sameMPV}
    \lean{MPSTensor.forall_k_exists_j_nondecaying_overlap_of_sameMPV}
    \leanok
    \uses{def:sector_bnt_cf, def:same_mpv2_pos,
        thm:sector_bnt_proportional_full_basis_match}
    Let $P, Q$ be two BNT canonical forms with the same MPV family at every
    positive length. Then for every sector $k$ of $Q$ there exist a sector $j$
    of $P$ with matched bond dimension $\dim_P(j) = \dim_Q(k)$, a gauge-phase
    equivalence after identifying the matched bond dimensions, and a
    non-decaying cross-overlap.
\end{theorem}

\begin{proof}\leanok
    \uses{thm:sector_bnt_proportional_full_basis_match}
    Convert equality to eventual nonzero proportionality with scalar $1$ and
    apply Theorem~\ref{thm:sector_bnt_proportional_full_basis_match}.
\end{proof}

\subsection{Bijective matching by symmetry}
\label{subsec:sector_bnt_strong_match_bijective_symmetry}

Applying the exact matching once more with $P$ and $Q$ swapped gives the reverse
direction; the basis-distinctness clause forces injectivity both ways, so
$g_a = g_b$ and the source relabelling
$B_k = e^{i\phi_k} X_k A_{j_k} X_k^{-1}$ holds on the full basis
\cite[Appendix~A]{Cirac2016MPDO_arXiv}.

\begin{theorem}[Bijective matching by symmetry]
    \label{thm:sector_bnt_bijective_match_of_sameMPV}
    \lean{MPSTensor.bijective_match_of_sameMPV}
    \leanok
    \uses{thm:sector_bnt_forall_unit_k_exists_j_nondecaying_overlap_of_sameMPV,
        lem:sector_bnt_bijection_from_matches,
        def:sector_bnt_cf, def:same_mpv2_pos}
    Let $P, Q$ be two BNT canonical forms with the same MPV family at every
    positive length. Then there is a bijection
    $\beta : J(Q) \simeq J(P)$ such that for every $k \in J(Q)$ the bond
    dimensions match, the basis tensors $P_{\beta(k)}$ and $Q_k$ are
    gauge-phase equivalent after identifying the matched bond dimensions, and
    the basis cross-overlap does not tend to zero.
\end{theorem}

\begin{proof}\leanok
    \uses{thm:sector_bnt_forall_unit_k_exists_j_nondecaying_overlap_of_sameMPV,
        lem:sector_bnt_bijection_from_matches}
    Apply the full-basis matcher in both directions, using symmetry of the
    same-MPV relation. Lemma~\ref{lem:sector_bnt_bijection_from_matches}
    converts these two families of matches into the required bijection.
\end{proof}
